\documentclass[aspectratio=169]{beamer}
\usetheme{Madrid}
\usepackage[utf8]{inputenc}
\usepackage[spanish]{babel}
\usepackage{amsmath, amssymb}
\usepackage{booktabs}
\usepackage{hyperref}

\title{Laboratorio interactivo de planificación de expansión de generación}
\subtitle{Caso Garver con AMPL e IA}
\author{Ing. Wilian Guamán Cuenca, M.Sc.}
\institute{ESPOCH -- Electricidad}
\date{2026}

\begin{document}

\begin{frame}
  \titlepage
\end{frame}

\begin{frame}{Propósito de la sesión}
  \begin{itemize}
    \item Usar un caso de estudio para comprender un modelo GEP.
    \item Explorar datos de tecnologías, demanda, reserva y ENS.
    \item Interpretar resultados de AMPL.
    \item Auditar explicaciones generadas por IA.
  \end{itemize}
\end{frame}

\begin{frame}{Repositorio docente}
  \begin{block}{Repositorio}
  \url{https://github.com/ESPOCH-Electricidad/power-systems-planning-operation}
  \end{block}
  \begin{itemize}
    \item Casos de estudio.
    \item Guías de modelos matemáticos.
    \item Notebooks de exploración.
    \item Actividades y rúbricas.
  \end{itemize}
\end{frame}

\begin{frame}{Organización por bloques}
  \begin{enumerate}
    \item Fundamentos de optimización.
    \item Operación de corto plazo.
    \item Flujo óptimo de potencia.
    \item Expansión de transmisión.
    \item Expansión de generación.
    \item Casos de estudio, evaluación e IA.
  \end{enumerate}
\end{frame}

\begin{frame}{Actividad inicial}
  \textbf{Pregunta Mentimeter:}

  En un modelo GEP, ¿qué parámetro condiciona más la decisión de expansión?

  \begin{itemize}
    \item Demanda.
    \item Inversión.
    \item Costo operativo.
    \item Reserva.
    \item ENS.
    \item Emisiones.
  \end{itemize}
\end{frame}

\begin{frame}{Modelo GEP: idea general}
  \begin{equation*}
  \min \sum_t\sum_g C_g P_{g,t} + \sum_t\sum_c IC_c I_{c,t} + \sum_t C^{ENS} ENS_t
  \end{equation*}

  El modelo decide capacidad nueva, despacho y energía no servida para atender demanda futura al menor costo.
\end{frame}

\begin{frame}{Restricciones principales}
  \begin{align*}
  Cap_{g,t} &= P_g^0 + \sum_{\tau\leq t} I_{g,\tau} \\
  \sum_g P_{g,t}+ENS_t &= D_t \\
  P_{g,t} &\leq Cap_{g,t} \\
  \sum_g firm_g Cap_{g,t} &\geq (1+RM)D_t
  \end{align*}
\end{frame}

\begin{frame}{Caso Garver}
  \begin{itemize}
    \item Tecnologías existentes: gas, hidro, térmica.
    \item Tecnologías candidatas: solar, eólica, ciclo combinado, hidro.
    \item Periodos: 2025, 2030, 2035.
    \item Restricciones: demanda, reserva, emisiones y ENS.
  \end{itemize}
\end{frame}

\begin{frame}{Notebook de exploración}
  \begin{itemize}
    \item Lectura de datos del caso.
    \item Gráficos de demanda.
    \item Comparación de costos, emisiones y capacidad firme.
    \item Verificación previa antes de resolver AMPL.
  \end{itemize}
\end{frame}

\begin{frame}{Actividad colaborativa}
  En Padlet o Google Docs, responder:
  \begin{enumerate}
    \item ¿Qué tecnología parece atractiva y por qué?
    \item ¿Qué pasa si aumenta la demanda?
    \item ¿Qué pasa si baja la disponibilidad hidroeléctrica?
    \item ¿Qué error podría cometer una IA al interpretar el caso?
  \end{enumerate}
\end{frame}

\begin{frame}{Desconfía del chatbot}
  \begin{quote}
  ``Como ENS es cero, el sistema es completamente seguro y no requiere análisis de red.''
  \end{quote}

  \textbf{Tarea:} detectar el error, corregirlo y justificarlo técnicamente.
\end{frame}

\begin{frame}{Evaluación final}
  \begin{itemize}
    \item Quiz en Google Forms, Kahoot o Quizizz.
    \item Preguntas sobre capacidad, despacho, reserva, ENS y uso crítico de IA.
    \item Reflexión final: ¿qué conclusión de IA fue más peligrosa?
  \end{itemize}
\end{frame}

\begin{frame}{Cierre}
  La misma lógica de laboratorio puede extenderse a ED, UC, OPF, TNEP y GEP multianual.
\end{frame}

\end{document}
